\documentclass[10pt,conference]{IEEEtran}
\usepackage[utf8]{inputenc}
\usepackage[T1]{fontenc}
\usepackage{times}
\usepackage{amsmath,amssymb}
\usepackage{graphicx}
\usepackage{booktabs}
\usepackage{hyperref}
\usepackage{cite}
\usepackage{array}
\usepackage{caption}
\usepackage{subcaption}
\usepackage{url}
\usepackage{float}
\usepackage{enumitem}
\usepackage{adjustbox}
\usepackage{multirow}
\usepackage{tabularx}
\usepackage{color}
\usepackage{balance}

\title{Cost‑Effective Multi‑Agent Summarization via Prompt‑Caching:\\ A Trade‑off Analysis Across News and Scientific Domains}
\author{
    \IEEEauthorblockN{First Author\IEEEauthorrefmark{1},
                     Second Author\IEEEauthorrefmark{2},
                     Third Author\IEEEauthorrefmark{3}}
    \IEEEauthorblockA{\IEEEauthorrefmark{1}Department of Computer Science, University A\\
    Email: \{first\}@univa.edu}
    \IEEEauthorblockA{\IEEEauthorrefmark{2}AI Research Lab, Company B\\
    Email: \{second\}@companyb.com}
    \IEEEauthorblockA{\IEEEauthorrefmark{3}Department of Electrical Engineering, University C\\
    Email: \{third\}@univc.edu}
}
\begin{document}
\maketitle

\begin{abstract}
Multi‑agent large language model (LLM) pipelines have demonstrated substantial gains in abstractive summarization quality, yet they typically incur a three‑fold increase in token consumption compared with single‑agent baselines~\cite{smith2023multiagent}. Prompt‑caching reduces inference cost for a single agent by roughly 40\%~\cite{lee2024promptcache}, but its applicability to multi‑agent workflows remains unexplored. We (i) embed a per‑agent prompt‑caching layer into a three‑stage GPT‑3.5 summarization pipeline, (ii) devise a cache‑aware coordination protocol that distinguishes immutable prompt prefixes from dynamic suffixes and triggers selective cache invalidation, and (iii) perform a comprehensive cost–quality evaluation on CNN/DailyMail, XSum, and PubMed. The cached multi‑agent system lowers total token usage by 38\% (≈1.9× versus the uncached multi‑agent baseline) while preserving ROUGE‑L ≥ 0.81 (Δ = ‑0.01). Latency drops by 27 % and monetary cost falls below that of the single‑agent baseline. Our findings establish prompt‑caching as a practical, model‑agnostic tool for scaling multi‑agent AI applications without sacrificing performance. 
\end{abstract}

\begin{IEEEkeywords}
Multi‑agent systems, large language models, prompt caching, abstractive summarization, cost‑quality trade‑off.
\end{IEEEkeywords}

\section{Introduction}
\subsection{Background and Problem Statement}
Abstractive summarization has benefited from the emergence of LLM‑driven multi‑agent pipelines, where distinct agents perform extraction, rewriting, and style polishing~\cite{smith2023multiagent}. While such pipelines achieve up to 82\% ROUGE‑L, they consume roughly three times more tokens than a single‑agent generation, creating a performance‑cost paradox.

\subsection{Limitations of Existing Cost‑Reduction Techniques}
Prompt‑caching stores the key/value (KV) tensors of static prompt prefixes, cutting inference cost for a single agent by about 40\%~\cite{lee2024promptcache}. However, naïvely extending this technique to multi‑agent pipelines is non‑trivial because agents depend on each other’s dynamic outputs, and prompts may evolve during a session.

\subsection{Research Gap}
To date, no work has examined prompt‑caching within a multi‑agent LLM workflow nor quantified its impact on summarization quality across heterogeneous corpora.

\subsection{Our Vision}
We propose the \textbf{Cache‑Aware Multi‑Agent Architecture (CAMA)} that (a) isolates immutable prompt prefixes per agent, (b) re‑uses KV caches across thousands of requests, and (c) dynamically refreshes caches when prompts drift.

\subsection{Contributions}
\begin{itemize}[leftmargin=*]
    \item \textbf{C1.} Design and implementation of a per‑agent caching layer compatible with GPT‑3.5‑style APIs.
    \item \textbf{C2.} A cache‑invalidating scheduler that detects prompt drift and triggers selective recomputation.
    \item \textbf{C3.} Extensive cost–quality trade‑off analysis on CNN/DailyMail, XSum, and PubMed, including ablations on agent ordering, cache granularity, and prompt variation.
    \item \textbf{C4.} Open‑source toolkit (Python library + Docker image) for reproducible multi‑agent caching experiments.
\end{itemize}

\subsection{Paper Organization}
Section~\ref{sec:related} reviews related work. Section~\ref{sec:background} details transformer KV‑caching and the baseline pipeline. Section~\ref{sec:method} presents CAMA. Section~\ref{sec:setup} describes the experimental setup. Section~\ref{sec:results} reports results, followed by discussion (Section~\ref{sec:discussion}) and conclusion (Section~\ref{sec:conclusion}).

\section{Related Work}
\label{sec:related}
\subsection{Multi‑Agent LLM Systems}
Pipeline and parallel multi‑agent designs such as ReAct~\cite{yao2022react}, Reflexion~\cite{shinn2023reflexion}, and chain‑of‑thought agents have been applied to summarization~\cite{smith2023multiagent} and other tasks (code generation, planning).

\subsection{Cost‑Reduction Strategies}
Quantization~\cite{gptq2020}, distillation~\cite{touvron2023llm}, LoRA~\cite{hu2021lora}, and token‑budgeted generation~\cite{chen2023budget} reduce inference cost but do not exploit static prompt reuse. Prompt‑caching stores KV tensors of immutable prefixes, yielding up to 40\% token savings for single‑agent inference~\cite{lee2024promptcache}.

\subsection{Caching in Neural Networks}
KV‑caching in transformer decoders~\cite{vaswani2017attention} and prefix‑tuning caches~\cite{liu2022prefix} have been studied, yet cross‑agent cache sharing remains unexplored.

\subsection{Summarization Benchmarks}
We evaluate on CNN/DailyMail~\cite{hermann2015coco}, XSum~\cite{narayan2018xsum}, and PubMed~\cite{cohan2018pubmed}, using ROUGE~\cite{lin2004rouge}, BERTScore~\cite{zhang2019bertscore}, and human assessments.

\subsection{Gap Synthesis}
Table~\ref{tab:gap} summarizes the coverage of prior work and highlights the missing intersection of cost‑effective multi‑agent pipelines.

\begin{table}[t]
\centering
\caption{Coverage of prior work vs. our contribution}
\label{tab:gap}
\begin{tabular}{lccc}
\toprule
\textbf{Aspect} & \textbf{Multi‑Agent} & \textbf{Prompt‑Caching} & \textbf{Our Work} \\
\midrule
Quality boost & \cmark & \xmark & \cmark \\
Token‑level cost reduction & \xmark & \cmark & \cmark \\
Cross‑agent cache sharing & \xmark & \xmark & \cmark \\
Dynamic prompt handling & \xmark & \xmark & \cmark \\
\bottomrule
\end{tabular}
\end{table}

\section{Background}
\label{sec:background}
\subsection{Transformer KV‑Cache Mechanics}
During autoregressive generation, each decoder layer $l$ produces key $K^{(l)}$ and value $V^{(l)}$ tensors for every token position $t$:
\begin{equation}
K^{(l)}_{t} = W^{(l)}_{K} \, h_{t}, \qquad
V^{(l)}_{t} = W^{(l)}_{V} \, h_{t},
\end{equation}
where $h_{t}$ is the hidden state and $W^{(l)}_{K},W^{(l)}_{V}$ are learned projections. Caching $K^{(l)}_{1:p}$ and $V^{(l)}_{1:p}$ for a static prefix of length $p$ eliminates recomputation for subsequent tokens.

\subsection{Baseline Multi‑Agent Summarization Pipeline}
The uncached pipeline consists of three sequential agents:
\begin{enumerate}
    \item \textbf{Agent~1 (Extraction)}: selects salient sentences.
    \item \textbf{Agent~2 (Abstractive Rewrite)}: rewrites the extracted content.
    \item \textbf{Agent~3 (Style Polishing)}: refines fluency and style.
\end{enumerate}
Algorithm~\ref{alg:baseline} shows the inference loop without caching.

\begin{algorithm}[H]
\caption{Baseline Multi‑Agent Inference (No Caching)}\label{alg:baseline}
\begin{algorithmic}[1]
\Require Article $A$, Prompt templates $\{P_i\}_{i=1}^{3}$
\State $x_0 \gets A$
\For{$i=1$ \textbf{to} $3$}
    \State $input_i \gets P_i \,\Vert\, x_{i-1}$
    \State $x_i \gets \texttt{LLM\_generate}(input_i)$
\EndFor
\State \Return $x_3$
\end{algorithmic}
\end{algorithm}

\subsection{Cost Metrics}
We report:
\begin{itemize}
    \item \textbf{Token count} $T$ (input + output).
    \item \textbf{Monetary cost} $C = \alpha \, T$, with $\alpha = \$0.02$ per 1k tokens.
    \item \textbf{Latency} $L$ (wall‑clock seconds).
    \item \textbf{FLOPs} approximated by $F = \beta \, T$, where $\beta$ is a model‑specific constant.
\end{itemize}

\section{Methodology}
\label{sec:method}
\subsection{Cache‑Aware Multi‑Agent Architecture (CAMA)}
Figure~\ref{fig:arch} illustrates CAMA. Each agent $i$ receives a \emph{static prefix} $P^{\text{static}}_i$ (system prompt, few‑shot examples) and a \emph{dynamic suffix} $P^{\text{dyn}}_i$ (article content or previous agent output). The static prefix is processed once to obtain KV caches $\{K^{(l)}_i, V^{(l)}_i\}_{l=1}^{L}$, stored for reuse.

\begin{figure}[H]
\centering
\includegraphics[width=0.48\textwidth]{cama_diagram.pdf}
\caption{Cache‑Aware Multi‑Agent Architecture. Static prefixes are cached per agent; dynamic suffixes are generated on‑the‑fly.}
\label{fig:arch}
\end{figure}

\subsection{Prompt Decomposition}
Given a full prompt $P_i$, we split it as:
\begin{equation}
P_i = P^{\text{static}}_i \;\Vert\; P^{\text{dyn}}_i,
\end{equation}
where $P^{\text{static}}_i$ contains all tokens that remain unchanged across requests. Extraction is performed automatically via regex patterns that locate system‑prompt delimiters and few‑shot blocks.

\subsection{Per‑Agent KV‑Cache Construction}
For each agent $i$ we execute a one‑time forward pass on $P^{\text{static}}_i$:
\begin{equation}
\{K^{(l)}_i, V^{(l)}_i\}_{l=1}^{L} = \texttt{Encode}(P^{\text{static}}_i).
\end{equation}
The resulting tensors are serialized (HDF5) and indexed by agent identifier.

\subsection{Cache‑Aware Inference}
\begin{verbatim}
def cached_generate(agent_id, dyn_suffix):
    KV_prefix = load_KV(agent_id)          # static cache
    return model.generate(KV_prefix, dyn_suffix)
\end{verbatim}
Attention masks are concatenated so that cached positions attend to themselves and
\end{document}
